\section{Reproducibility Runbook}
\label{app:repro}

This appendix is the step-by-step reproduction runbook. Every kernel
lives at the canonical hexa-lang stdlib home (\verb|stdlib/cern/|, @D d3)
and is callable from \verb|hexa verify --expr <fn> <args> <expected>|.
The Bethe--Bloch and plasma-wakefield kernels are hexa-native; the Xsuite
optics (xsuite\,0.51.1) and pylhe statistics (pylhe\,1.0.4) adapters and
the FBPIC parity run are Python under the same home. Artifacts land under
\verb|exports/cern/|.

% -------------------------------------------------------------------------
\subsection{Per-atom \texttt{hexa verify} invocations}
\label{app:repro-verify}

\begin{lstlisting}
# Bethe-Bloch dE/dx (Stage-3 GREEN, 4 PDG anchors)
hexa run stdlib/cern/bethe_bloch_stopping.hexa      # selftest 4/4 PASS
hexa verify --expr bethe_bloch_stopping "Al" 1.0    178.824652
hexa verify --expr bethe_bloch_stopping "Pb" 1000.0 1.196980424

# Plasma-wakefield cold-linear (selftest 5/5 GREEN)
hexa run stdlib/cern/plasma_wakefield.hexa
hexa verify --expr wakefield_e0_gv_m   1e18   96.15919872735485
hexa verify --expr wakefield_lambda_um 1e18   33.38943270215429

# 1-D linear PIC parity (FBPIC pins; pic_parity_check)
hexa run stdlib/cern/plasma_wakefield.hexa --pic    # Delta=3.562% PASS

# Xsuite FODO twiss (Python adapter, parity oracle built in)
python3 stdlib/cern/xsuite_optics.py                # beta/q rel ~1e-14

# pylhe LHE round-trip (Python adapter; absorbed=false)
python3 stdlib/cern/lhe_stats.py                    # GATE_OPEN (permanent)
\end{lstlisting}

% -------------------------------------------------------------------------
\subsection{Adapter environment pins}
\label{app:repro-env}

\begin{table}[h]
\centering
\small
\begin{tabular}{@{}l l l@{}}
\toprule
component & version & host / note \\
\midrule
hexa (kernels)   & strict mode      & local build \\
Xsuite           & 0.51.1           & Python 3.x, optics-only \\
pylhe            & 1.0.4 (BSD-3)    & Python 3.14.4 \\
FBPIC            & 0.27.0           & pool \texttt{ubu-1}, CPU/numba \\
\texttt{particle}& 0.26.2           & Stage-1 substrate ref capture \\
\bottomrule
\end{tabular}
\caption{Adapter / backend version pins. The FBPIC parity ran on a free
Linux pool host (no paid compute); everything else is local.}
\label{tab:app-repro-env}
\end{table}

% -------------------------------------------------------------------------
\subsection{Artifact tree and the build}
\label{app:repro-tree}

\begin{lstlisting}
exports/cern/
  stopping/<stamp>/   cern_g4_stopping_*.json   # 28-row dE/dx sweep
  synthesize/<stamp>/ cern_synth_*.json         # FODO twiss + parity
  lhe/<stamp>/        cern_lhe_*.json           # LHE round-trip stats
  {specify,structure,design,handoff}/           # stub dirs (rc=2, empty)

cern-accelerator/         # this monograph
  main.tex  references.bib  Makefile
  appendix/A_*.tex ... J_*.tex
  figures/_scripts/figNN_*.tex  ->  figures/figNN_*.pdf
  companion/verify-ledger.json  pr-roll.json  session-journal.md
\end{lstlisting}

% -------------------------------------------------------------------------
\subsection{Bit-for-bit fingerprint}
\label{app:repro-fp}

The stopping-power export carries a content fingerprint that is stable
across every re-run --- the strongest possible reproduction signal. Ten
independent runs of the dE/dx sweep over two days
(\texttt{2026-05-19}--\texttt{2026-05-21}) all emit the same digest:

\begin{lstlisting}
$ jq -r '.fingerprint' exports/cern/stopping/*/cern_g4_stopping_*.json | sort -u
60e34a15cfb4477f      # one line -> all 10 runs bit-identical
\end{lstlisting}

Because the kernel is a deterministic closed-form evaluation (no RNG, no
iterative solver) the fingerprint is a function of the inputs alone, so a
reader who reruns the sweep and gets \texttt{60e34a15cfb4477f} has
reproduced the data record exactly. The LHE round-trip is likewise
deterministic (RNG seed \texttt{20260520} pinned in the record), and the
plasma-wakefield / FODO atoms are pure closed-form recomputes.

\paragraph{End-to-end walkthrough.} Clone the two repos; run the
\verb|hexa verify| / \verb|hexa run| invocations above to reproduce the
ledger values; then build the monograph. The build is \texttt{make}
(xelatex $\times 3$ + bibtex); figures are regenerated from their
\texttt{.tex} sources with \texttt{make figures} (each standalone
pgfplots snippet compiles to a \texttt{figures/figNN\_*.pdf}); page and
word counts with \texttt{make pages} / \texttt{make wordcount}; and lint
with \texttt{make lint} (commons @D g51 --- pages + fal.ai figure floor).
The six pgfplots charts (Bethe--Bloch dE/dx, wakefield density scaling,
PIC parity, FODO $\beta$-profile, tier distribution) plus the fal.ai
cover and the TikZ pipeline diagram are all reproducible from the sources
under \verb|figures/|.
